\section{The typed maps to the other programmes}\label{sec:bridges}

The workbench connects the Navier--Stokes construction to the author's other programmes through explicit maps. The table states what each map carries and what is established about it here; \emph{verified} marks an identity or a statement re-derived here.

\begin{center}\small
\begin{longtable}{@{}>{\raggedright\arraybackslash}p{0.2\linewidth}>{\raggedright\arraybackslash}p{0.36\linewidth}>{\raggedright\arraybackslash}p{0.36\linewidth}@{}}
\toprule
Map & What it carries & Established here\\
\midrule
\endhead
NS $\to$ YM, abelian & $A_i=\lambda u_iT$; the curvature masses $\mathcal K_B=\lambda^2\|\operatorname{curl}u\|^2$ and $\mathcal K_E=\lambda^2c^{-2}\|\partial_tu\|^2$ & the map and its identities, verified; the divergence of the curvature masses conditional on (H1)--(H4) (Theorem~\ref{thm:rates}); the zero-quotient diagonal needs couplings $g_j\to0$\\
NS $\to$ YM, non-abelian & $\mathbf A_i=\lambda e_i\times u$; its fields, invariants and currents & all identities verified for every divergence-free $u$; the off-axis vorticity bound of Proposition~\ref{prop:typeII}(b) under (H1)--(H4). The exact axis identity of the source's (243) comes from the workbench's own axis continuation, not from (H1)--(H4), since (5.42) does not reach the axis. The current is nonzero, so the image is a sourced field\\
NS $\to$ Erdős--Straus & the auxiliary torus operator $J=\begin{psmallmatrix}3&1\\1&5\end{psmallmatrix}$, with $\det J=14$ and eigenvalues $4\pm\sqrt2$ & the operator identities used in $21$ statements of the Erdős--Straus supplement, verified in the Erdős--Straus paper \cite{ESPaper}\\
NS $\to$ zeta & angular-momentum balance, a Mellin shift and the Müntz multiplier $W_h(s)=\zeta(s)(1-2^{s-1})(1-2^{h-s})$ & the transforms, verified; the growth, conditional on Theorem~\ref{thm:rates}; the Mellin abscissa of $\mathcal K_B(1-\tau)$ lies in $[\frac12+3h,1]$, conditionally\\
$S^6\to$ NS & a Beltrami-mode forced solution on $\mathbb T^3$ from the $S^6$ oscillator & a global smooth solution, verified exactly\\
Jacobian map $\to$ NS kinematics & $U=DF^{-1}(V\circ F)$, a divergence-free polynomial field whose flow escapes at $\tau=1/8$ & verified exactly (the Jacobian paper \cite{JacobianPaper}, §4)\\
NS $\to$ Einstein & the snapshot encoding, the Kasner map, the Rindler response and the slab comparator & all identities verified; the collision and the exact radius of Section~\ref{sec:rindler}\\
\bottomrule
\end{longtable}
\end{center}

My assessment of these maps, taken together: the map to Einstein gravity has produced the strongest result so far, the exact radius of Section~\ref{sec:rindler}. I think it is of independent interest in fluid/gravity duality, because it places the breakdown of the hydrodynamic expansion of the Rindler fluid at a certified pole collision at real momentum, a question the continuation had left open. The maps to Yang--Mills are exact as identities, and their analytic content is conditional on (H1)--(H4); certifying those four statements (Question~\ref{q:profiles}) would make the curvature rates unconditional. For the maps to the zeta programme and to the Erdős--Straus workbench, what is established is the exact identities; what they carry about the zeta function or the equation beyond those identities is not verified here.

The connections to established work are these: the convergence of the hydrodynamic gradient expansion \cite{GKST,Withers,HSSSW}; the breakdown of diffusion by pole collision \cite{ADGS}; fluid/gravity duality and the Rindler fluid \cite{BHMR,BKLS,CMST,MarolfRangamani}; forced singularities, whose strategy of the force as a residual with controlled regularity the manuscript shares \cite{CordobaMartinezZoroa}; the classical blowup criteria of Section~\ref{sec:rates}; and singularity analysis of power series \cite{FlajoletOdlyzko}. The manuscript's Lemma~10.3 is, in its structure, Borel's lemma with shrinking cutoffs (a reading of the first pass), and its exterior heat profile $\mathcal H(Z)=Z^{-1-h}U(1+h,2,1/Z)$ is a Tricomi function (verified in \note{first\_pass/check\_axis\_numerics.py}).
